\section{An obstruction to Booleanizing the canonical point}
\label{sec:booleanization}

The distinction between an extension somewhere in the class and an
extension inside a fixed model also affects the topos argument in
\cite[Theorem~4.1]{Espindola2023Spectra}. We give a concrete obstruction
on the site of countable models: adding a base point is always possible
there, but need not be possible inside a given countable-base model.
The resulting dense cosieve prevents that model's canonical presheaf
point from factoring through the full double-negation subtopos.
Throughout this section we retain CH and the class constructed above.

\subsection{Adding base points to a countable model}

\begin{lemma}\label{bool:extension}
Every countable \(M\in K\) has a countable strong extension \(M'\)
with \(I^{M'}=I^M\) and \(B^M\subsetneq B^{M'}\).
\end{lemma}

\begin{proof}
First we construct compatible colors for all tuples over \(I=I^M\).
Choose an injection \(j:I\to\omega\). For each nonempty finite
\(X\subseteq I\), choose a surjection \(d_X:\omega\to\omega\),
requiring \(d_{\{x\}}(0)=j(x)\). For example, set
\(d_X(n+1)=n\), and choose the zeroth value as prescribed for
singletons and as zero otherwise. For a tuple \(t=(x_1,\ldots,x_k)\)
and a label \((J,n)\in\mathcal L_k\), put
\[
 q_t(J,n)=d_{\{x_r:r\in J\}}(n).
\]
Define a diagram \(D_t\) by equality of these values, by ordinal
membership between them, and by the tag \(\rho_t(J,n)=q_t(J,n)\).
Its quotient is \(\omega\), since the labels \(([k],n)\) already
surject onto \(\omega\). Ordinal membership on this quotient is
extensional and strictly decreases the tags, which lie below
\(\delta\). The distinguished singleton labels have the distinct
values \(j(x_r)\). Thus \(D_t\) is valid. Restricting to a
subsequence gives exactly its separately valid diagram: the definition
of \(d_X\) depends on the underlying finite set, not its position in
a larger tuple. Consequently the colors
\(s_t=\iota_k(D_t)\) satisfy every full-subsequence test. When \(I\)
is infinite, the unary colors are distinct; when \(I\) is finite,
there are no tuples at sufficiently large arities. Thus this does not
produce the fixed-arity pattern excluded by Lemma~\ref{col:no-pattern}.

Take the ordered sum \(B'=B^M+\Q\), with a disjoint copy of \(\Q\).
It is a countable dense order without endpoints, and \(B^M\) is
convex in it. Keep \(I^M\) and the exact old labeled group \(G_*^M\).
Enlarge \(G(B^M)\) to \(G(B')\) by retaining each old code with
its old support and adding one new code for every remaining finite
subset of \(B'\). Choose old row origins \(z_{t,n}\), and extend
each old row torsor to a \(G(B')\)-torsor: retain the points with
old support coordinates and adjoin one point for each remaining
coordinate, with the action given by symmetric difference.
Keep every old column torsor, choose an origin
\(w_{t,b}\) in it, and add a new \(G_*^M\)-torsor with an origin
\(w_{t,b}\) for each new \(b\).

Choose \(a_t\in2^\omega\) representing \(s_t\), and prescribe the
origin bits by
\[
 a_{t,b}(n)=
 \begin{cases}
 e^M(z_{t,n},w_{t,b}),&b\in B^M,\\
 a_t(n),&b\in B'\setminus B^M.
 \end{cases}
\]
Extend them to all points by the torsor formula
\[
 e'(z_{t,n}+u,w_{t,b}+v)
 =a_{t,b}(n)+\boldsymbol1_{b\in\supp(u)}
             +\boldsymbol1_{n\in v}\pmod2.
\]
Unique torsor coordinates make this well-defined. The old translation
law~\eqref{aec:bit} shows that every old bit, true or false, is
unchanged. All group and action values on old arguments are unchanged
as well. Give the enlarged carriers exactly the required relations,
with zero bits between different tuple fibers. This is a basic
inclusion, and its countable base makes the enlarged structure an
object.

At every new base point the colors, using the chosen old rows, are
the \(s_t\), so all old-tuple tests pass. Changing to other old row
points has no effect there: the translators have supports contained
in \(B^M\). Changing column origins preserves colors modulo
\(=^*\). Hence the inclusion is strong, not just basic. Finally,
Lemma~\ref{aec:size} makes the extension countable.
\end{proof}

\subsection{The site and the canonical point}

Let \(\mathcal C_0\) be a small skeleton of the category of countable
models in \(K\) and strong embeddings. We use covariant presheaves
\(\mathcal P=\mathrm{Set}^{\mathcal C_0}\), that is, presheaves on
\(\mathcal C_0^{\mathrm{op}}\), and write
\(y_A=\operatorname{Hom}_{\mathcal C_0}(A,-)\). A cosieve on \(A\)
is a collection of maps out of \(A\) closed under postcomposition,
or equivalently a subfunctor of \(y_A\). The double-negation topology
\(J_{\neg\neg}\) on \(\mathcal C_0^{\mathrm{op}}\) declares a
cosieve \(R\) covering exactly when
\begin{equation}\label{bool:density}
 \text{for every } f:A\to D\text{ there is }g:D\to E
 \text{ with }gf\in R.
\end{equation}
Thus \(\operatorname{Sh}(\mathcal C_0^{\mathrm{op}},J_{\neg\neg})\)
is the full double-negation Booleanization of \(\mathcal P\).

For a regular infinite cardinal \(\theta\), a poset is
\(\theta\)-directed if every subset of cardinality less than
\(\theta\) has an upper bound. A model \(M\) is
\(\theta\)-presentable if \(\operatorname{Hom}_K(M,-)\) preserves
\(\theta\)-directed colimits; its presentability rank is the least
regular infinite \(\theta\) with this property.
Our site is also the countable \emph{internal-size} site in the convention
of \cite[Definition~2.3 and the notation at the end of
Section~2]{Espindola2022Eventual}. Indeed, for every regular
\(\theta>\aleph_0\), a model of cardinality less than \(\theta\)
is \(\theta\)-presentable: its image in a \(\theta\)-directed
colimit lies in one stage, coherence makes the factorization strong,
and embeddings give uniqueness. Conversely, every model is the union
of its strong submodels of cardinality less than \(\theta\), a
\(\theta\)-directed family by the downward
L\"owenheim--Skolem property and regularity of \(\theta\).
If the model has cardinality at least \(\theta\), its identity
cannot factor through any member of this family. An uncountable model
of cardinality \(\nu\) therefore has presentability rank \(\nu^+\).
A countable model has rank either \(\aleph_0\) or \(\aleph_1\).
The cited convention assigns internal size \(\nu\) to rank
\(\nu^+\), and assigns a limit rank to itself. Thus internal and
underlying cardinalities agree here, including at \(\aleph_0\).

For any \(H\in K\), let \(\mathcal C_0/H\) have as objects the
strong maps \(c:C\to H\) with \(C\in\mathcal C_0\), and as arrows
the commuting triangles. It is \(\aleph_1\)-filtered: countably many
images are contained in a countable strong submodel of \(H\), and
coherence gives the required transition maps. Parallel arrows over
\(H\) are equal because the maps into \(H\) are injective.
The canonical point \(p_H:\mathrm{Set}\to\mathcal P\) has inverse
image
\begin{equation}\label{bool:point}
 p_H^*(X)=\mathop{\mathrm{colim}}_{(C,c)\in\mathcal C_0/H}X(C),
 \qquad p_H^*(y_A)=\operatorname{Hom}_K(A,H).
\end{equation}
The first formula preserves colimits and finite limits, giving a
point; in fact it preserves countable limits since the indexing
category is \(\aleph_1\)-filtered. In the second formula, a
representative \(r:A\to C\) maps to \(cr\). Every map into \(H\)
occurs this way, and two representatives with the same composite
agree after passing to a common countable strong submodel of \(H\).
No identification of all points of \(\mathcal P\) with models of
\(K\) is needed.

\begin{proposition}\label{bool:obstruction}
If \(H\in K\) has countable base, its canonical point \(p_H\)
does not factor through
\(\operatorname{Sh}(\mathcal C_0^{\mathrm{op}},J_{\neg\neg})\).
This applies, in particular, to every model on the categoricity tail
\(\|H\|\ge\Lambda\).
\end{proposition}

\begin{proof}
Choose a countable strong map \(a:A\to H\) whose image contains
\(B^H\), using the downward L\"owenheim--Skolem property and then
the representative in \(\mathcal C_0\). Define
\[
 R_A=\{r:A\to D\text{ in }\mathcal C_0:
                 r[B^A]\subsetneq B^D\}\subseteq y_A.
\]
Injectivity makes this a cosieve. For any \(f:A\to D\),
Lemma~\ref{bool:extension}, followed by an isomorphism to the chosen
representative in \(\mathcal C_0\), gives a countable strong map
\(g:D\to D'\) adding a base point outside \(g[B^D]\).
It follows that \(gf\in R_A\). Thus \(R_A\) satisfies
\eqref{bool:density}, without any amalgamation assumption.

The image of \(p_H^*(R_A)\to p_H^*(y_A)\) consists exactly of
maps \(A\to H\) factoring as \(hr\) with \(r\in R_A\) and
\(h:D\to H\) strong. But \(a\) has no such factorization.
If \(a=hr\), choose \(b\in B^D\setminus r[B^A]\).
Since \(h(b)\in B^H=a[B^A]\), there is \(x\in B^A\) with
\(h(b)=a(x)=h(r(x))\), contrary to injectivity of \(h\).

Sheafification for \(J_{\neg\neg}\) inverts the dense inclusion
\(R_A\hookrightarrow y_A\). If \(p_H\) factored through the
sheaf subtopos, its inverse image would therefore send that inclusion
to an isomorphism. The omitted element \(a\) proves that it does
not even send it to a surjection. Finally, models of cardinality at
least \(\Lambda\) have countable base by
Lemmas~\ref{aec:order-size} and~\ref{aec:size} together with
Proposition~\ref{cat:bounded}, and exist by Theorem~\ref{cat:tail}.
\end{proof}

\subsection{The precise comparison}

The proof of \cite[Theorem~4.1]{Espindola2023Spectra}, in the cited
version, starts with the full double-negation subtopos on the
internal-size site and asserts that the large categorical model gives
a point of it. Proposition~\ref{bool:obstruction} rules out the
canonical factorization for the present class, even though its point
preserves countable limits. A topology generated by a smaller chosen
family of covers could omit \(R_A\); in that case it would not be
the full double-negation topology used in that step.

This is an explicit obstruction to that step, not a classification of
all points of the Booleanization or a proof that it has no points.
It does not analyze alternative constructions of points or establish
what additional hypotheses would suffice to repair the cited argument.
